\documentclass[11pt,letterpaper]{article}
\usepackage[utf8]{inputenc}
\usepackage{geometry}
\usepackage{setspace}

\geometry{left=1.0in,right=1.0in,top=1.0in,bottom=1.0in}
\doublespacing

\setlength{\parindent}{0pt}
\setlength{\parskip}{12pt plus 2pt minus 2pt}

\begin{document}

\thispagestyle{empty}

\begin{center}
    \large\textbf{IN THE UNITED STATES PATENT AND TRADEMARK OFFICE} \\
\end{center}

\vspace{0.2cm}
\noindent\begin{tabular}{p{2.3in}p{4.0in}}
\textbf{In re Patent Application of:} & Reinaldo Javier Aguiar Marcano \\
\textbf{Application Number:} & 19/560,062 \\
\textbf{Confirmation Number:} & 9071 \\
\textbf{Filing Date:} & March 8, 2026 \\
\textbf{Group Art Unit:} & Pre-Exam / TC 2100 \\
\textbf{Examiner:} & Andrea Bennett \\
\textbf{Title of Invention:} & PORTABLE DATA CARTRIDGE WITH OPTICAL INTERFACE \\
\end{tabular}

\vspace{0.6cm}
\begin{center}
    \textbf{\large SUBSTITUTE SPECIFICATION (37 CFR \S{} 1.125)} \\
    \textbf{(WITH COMPLIANT LINE SPACING UNDER 37 CFR \S{} 1.52(b))}
\end{center}

\vspace{0.4cm}

\begin{center}
    \textbf{\large PORTABLE DATA CARTRIDGE WITH OPTICAL INTERFACE}
\end{center}

\section*{Cross-Reference to Related Applications}
This application claims the benefit of and priority to U.S. Patent Application No. 19/559,705, filed on March 6, 2026, the contents of which are hereby incorporated by reference in their entirety.

\section*{Background of the Invention}
As datasets in fields like AI, genomics, and national security grow into the petabyte and exabyte scale, transferring them over conventional telecommunication networks has become increasingly impractical due to bandwidth limitations and security concerns. The alternative, physically transporting storage media (a practice often called ``sneakernet''), is limited by the slow I/O speeds of traditional portable hard drives or storage arrays. The time taken to load and unload the data at either end often negates the benefits of physical transport. There is a need for a portable storage device that combines high storage capacity with an ultra-high-bandwidth interface to make physical data logistics truly efficient.

\section*{Summary of the Invention}
The present invention is a portable data cartridge with a high-bandwidth optical interface, colloquially referred to as a ``data brick,'' that is specifically designed to solve the I/O bottleneck in physical data transport. The cartridge is a self-contained unit comprising three key components within a durable, shock-resistant housing: a high-capacity non-volatile storage medium, an internal power source, and a pair of massively parallel optical interfaces (one for input, one for output).

The optical interfaces are consistent with the CVBNI architecture described in the parent application, No. 19/559,705. They enable the cartridge to ingest or offload its entire data payload at multi-terabit-per-second speeds. This reduces the ``loading dock'' time for a petabyte of data from many hours to mere minutes, making the physical transport of the cartridge via courier, drone, or vehicle an extremely fast and efficient method for bulk data transfer. The cartridge is the key enabling apparatus for a new ecosystem of secure physical data logistics.

\section*{Statement of License Grant to the United States Government}
The inventor, Reinaldo Javier Aguiar Marcano, hereby grants to the United States Federal Government and its military and defense agencies a perpetual, royalty-free, non-exclusive license to practice and use the invention described herein for any and all governmental, defense, and non-commercial purposes. This license explicitly excludes any and all state and local governments.

\section*{Brief Description of the Drawings}
FIG. 1 is an isometric view of the Portable Data Cartridge with Optical Interface, showing its ruggedized housing and optical I/O ports.

FIG. 2 is a block diagram illustrating the internal components of the data cartridge.

\section*{Detailed Description of the Invention}
Referring to FIG. 1, the apparatus is a portable data cartridge (100). It is enclosed in a ruggedized, shock-absorbent housing (101) suitable for physical transport. The housing includes at least two massively parallel optical ports: an input port (102) for receiving data and an output port (103) for transmitting data. These ports are protected by a retractable, sealed cover (104) to prevent contamination of the optical surfaces during transit. The housing may also include status indicators (105) for power and data activity.

FIG. 2 is a block diagram showing the internal components of the cartridge (100). The core components are a high-capacity, non-volatile storage medium (201), a storage controller (202), and an integrated power source (203), such as a high-capacity rechargeable battery. The storage controller (202) manages the flow of data to and from the optical interfaces. The input optical port (102) is coupled with a photon sensor array (204), while the output optical port (103) is coupled with a photon emitter array (205). This bidirectional design with separate input and output arrays allows for maximum data transfer speed and efficiency.

\vspace{0.8cm}
\begin{center}
    \textbf{\large INVENTOR STATEMENT UNDER 37 CFR \S{} 1.125(b)}
\end{center}

Pursuant to 37 CFR \S{} 1.125(b), the undersigned Applicant hereby certifies that this substitute specification includes no new matter.

\vspace{0.6cm}
\begin{singlespace}
\noindent Respectfully submitted, \par\vspace{1.0cm}
\underline{\textbf{/Reinaldo Javier Aguiar Marcano/}} \\
\textbf{Reinaldo Javier Aguiar Marcano} \\
Pro Se Applicant \\
Mailing Address: 2302 Britton Ridge Drive, Katy, Texas 77494 \\
Date: October 9, 2026
\end{singlespace}

\end{document}
